% !TEX program = xelatex
% !TEX spellcheck = zh-cn
% ============================================
% 李久阳 · 个人简历
% 编译环境：Overleaf — 务必选择 XeLaTeX 编译器（Menu > Compiler > XeLaTeX）
%            pdfLaTeX 会因中文支持缺失而报错（unisong30 / fandol）
% 创建：2026-07-13 · 修复：2026-07-13
% ============================================

\documentclass[11pt,a4paper]{ctexart}

% ---- 页面布局 ----
\usepackage[top=2cm, bottom=2cm, left=2.2cm, right=2.2cm]{geometry}
\usepackage{parskip}

% ---- 字体与颜色 ----
\usepackage{fontawesome5}
\usepackage{xcolor}
\usepackage{amssymb}
\definecolor{primary}{HTML}{2c3e50}
\definecolor{accent}{HTML}{2980b9}
\definecolor{gray}{HTML}{7f8c8d}
\definecolor{lightgray}{HTML}{bdc3c7}

% ---- 链接 ----
\usepackage[hidelinks]{hyperref}
\hypersetup{colorlinks=true, urlcolor=accent}

% ---- 列表 ----
\usepackage{enumitem}
\setlist{nosep, leftmargin=*, after=\vspace{0.3em}}
\setlist[itemize]{label={\color{accent}\scriptsize\(\blacktriangleright\)}}

% ---- 页眉页脚 ----
\usepackage{fancyhdr}
\pagestyle{fancy}
\fancyhf{}
\renewcommand{\headrulewidth}{0pt}
\cfoot{\color{lightgray}\small\faIcon{code} 此简历由 \LaTeX\ 生成 · 久阳 \& 渡 协作出品}

% ============================================
% 自定义命令
% ============================================

% 章节标题
\newcommand{\cvsection}[1]{
  \vspace{0.8em}
  \noindent
  {\color{primary}\large\textbf{#1}}
  \vspace{0.15em}
  \hrule
  \vspace{0.5em}
}

% 条目：标题 + 时间 + 描述
\newcommand{\cventry}[3]{
  \noindent
  \textbf{#1} \hfill \color{gray}\small #2
  \par\vspace{0.15em}
  \begin{itemize}
    #3
  \end{itemize}
}

% 标签/技能小方块
\newcommand{\tagbox}[1]{
  \colorbox{accent!12}{\color{accent}\small\sffamily #1}
}

% ============================================
\begin{document}

% ---- 头部 ----
\begin{center}
  {\color{primary}\huge\textbf{李久阳}} \\[0.3em]
  {\large 应用物理学本科生} \\[0.5em]
  {\color{gray}\small
    \faIcon{envelope} \href{mailto:jiuyang@example.com}{jiuyang@example.com} \quad
    \faIcon{github} \href{https://github.com/jiuyang}{github.com/jiuyang} \quad
    \faIcon{map-marker-alt} 中国·北京
  }
\end{center}

\vspace{-0.5em}

% ============================================
% 教育背景
% ============================================
\cvsection{\faIcon{graduation-cap} 教育背景}

\cventry{XX大学 · 应用物理学}{2023--2027（预计）}{
  \item 主修课程：量子力学、统计物理、电动力学、固体物理、数学物理方法
  \item 自学延伸：信息论、认知科学、决策论、复杂系统
  \item 将物理学中的建模方法论迁移至跨域分析：从第一性原理出发，追问系统的最小基元与生成规则
}

% ============================================
% 研究兴趣
% ============================================
\cvsection{\faIcon{search} 研究兴趣}

\begin{itemize}
  \item \textbf{信息论与社会系统}：将 Shannon 信息论框架推广至社会传播、注意力分配与认知熵，探究“信念更新幅度”（KL 散度）作为信息价值的统一度量
  \item \textbf{认知协作系统}：LLM 时代下人-AI 协作的架构设计——外置记忆、规则协议、索引系统如何构成比 prompt engineering 更鲁棒的认知基础设施
  \item \textbf{复杂系统建模}：从物理学出发，为社会学、经济学、认知科学建立可量化的形式模型（相变模型、场模型、OS 层级模型）
\end{itemize}

% ============================================
% 项目经验
% ============================================
\cvsection{\faIcon{tools} 项目经验}

\cventry{笔默 —— AI for Reading}{2026.03--至今}{
  \item 纸质阅读的“第二层智能”系统：AI 贴附在原有阅读行为链路上，不替代，只增强
  \item 底层认知模型基于预测加工（Predictive Processing）理论——“理解”来自“预测—修正”闭环
  \item 独立完成完整商业计划书（十三章），设计硬件形态（笔/耳机/台灯）、AI 交互范式、社区生态
  \item 核心设计哲学“最小侵入式增强”可跨域迁移至厨房 OS、驾驶 OS、健身 OS
}

\cventry{忆阻器时间展宽微弱电流测量}{2026}{
  \item 利用忆阻器非易失性电荷积分特性，将 pA 级微弱电流测量转化为阻值/时间测量
  \item 方案设计与可行性论证阶段
}

\cventry{个人主页}{2026.07}{
  \item 从零构建自我语境——不依赖社交平台的默认框架，定义自己的身份生成器
  \item 思路梳理阶段
}

% ============================================
% 技能
% ============================================
\cvsection{\faIcon{cogs} 技能}

\noindent
\tagbox{系统化分析} \quad
\tagbox{跨域建模} \quad
\tagbox{协议设计} \quad
\tagbox{信息论} \quad
\tagbox{认知科学} \quad
\tagbox{边界条件穷举}

\vspace{0.5em}

\noindent
\tagbox{\LaTeX\ (学习中)} \quad
\tagbox{Python} \quad
\tagbox{Markdown} \quad
\tagbox{Git} \quad
\tagbox{物理实验设计} \quad
\tagbox{数学推导}

% ============================================
% 方法论工具箱
% ============================================
\cvsection{\faIcon{layer-group} 方法论工具箱}

\noindent
\textbf{分析工具}：公理化驱动（追问最小基元）· 边界条件穷举（失败模式相图/反向复原法/减法原则/反向逼问法）· 比较矩阵（参数化→维空间对比）· 元层次追问（框架合法性而非内部正确性）

\vspace{0.3em}

\noindent
\textbf{建构工具}：协议化思维（将模糊判断转化为显式规则）· 迭代式精炼（提出→反馈→修正闭环）· 递归式自我应用（方法论折返应用于自身）· 反玄学五步法（语言还原/系统降维/举证反弹/直觉去魅/科学制度化）

% ============================================
% 语言
% ============================================
\cvsection{\faIcon{language} 语言}

\noindent
中文（母语） \hfill 英文（学术阅读与写作） \hfill 德语（入门）

\vspace{1.5em}
\begin{center}
  {\color{lightgray}\rule{0.5\textwidth}{0.3pt}} \\[0.3em]
  {\color{gray}\footnotesize \faIcon{quote-left} \ 我相信，人类的好奇心可以延向多远的地方。 \ \faIcon{quote-right}}
\end{center}

\end{document}
